% part: intuitionistic-logic
% chapter: propositions-as-types
% section: proofs-to-terms

\documentclass[../../../include/open-logic-section]{subfiles}

\begin{document}

\olfileid{int}{pty}{pt}

\olsection{Converting \usetoken{P}{derivation} to Proof Terms}

The process of converting natural deduction !!{proof}s to proof terms,
i.e., translating !!{proof}s in~\Log{N2i} to !!{proof}s in~\Log{tN2i},
is extremely straightforward. It simply amounts to writing a proof
term to the left of each !!{formula} on the right of a sequent in the
!!{derivation}, resulting in sequents of the form $\Gamma \Sequent M :
!A$.  We'll then say that $M$ \emph{witnesses}~$!A$ in
context~$\Gamma$. Which proof term to write is entirely determined by
the rules of \Log{tN2i}.

\begin{prop}
Every !!{proof}~$\delta$ of $\Gamma \Sequent !A$ in \Log{N2i}
corresponds to !!a{proof}~$\delta'$ of $\Gamma \Sequent N : !A$ in
\Log{tN2i}. The proof term $N$ is uniquely determined by~$\delta$.
\end{prop}

\begin{proof}
We proceed by induction on the height of~$\delta$. If $\delta$ is the
axiom $x:!A \Sequent !A$, then $\delta'$ is the axiom~$x:!A \Sequent
x:!A$. If $\delta$ ends in an inference rule, the inductive hypothesis
provides proof terms for each premise, and \Log{tN2i}-!!{proof}s of
sequents including these proof terms  corresponding to each premise.
We apply the corresponding rule of \Log{tN2i}; the corresponding proof
term is determined by the \Log{tN2i} inference rule.
\end{proof}

\begin{ex}
  Consider the very simple !!{proof} in \Log{N1i} of $(!A \land !B)
  \lif (!B \land !A)$:
  \[
  \AxiomC{$\Discharge{!A\land !B}{x}$}
  \RightLabel{\Elim\land[2]}
  \UnaryInfC{$!B$}
  \AxiomC{$\Discharge{!A\land !B}{x}$}
  \RightLabel{\Elim\land[1]}
  \UnaryInfC{$!A$}
  \RightLabel{\Intro\land}
  \BinaryInfC{$!B \land !A$}
  \DischargeRule{\Intro\lif}{x}
  \UnaryInfC{$(!A \land !B) \lif (!B \land !A)$}
  \DisplayProof
  \]
  In \Log{N2i} it looks like this:
  \[
  \Axiom$x : !A\land !B \fCenter !A\land !B$
  \RightLabel{\Elim\land[2]}
  \UnaryInf$x : !A\land !B \fCenter !B$
  \Axiom$x : !A\land !B \fCenter !A\land !B$
  \RightLabel{\Elim\land[1]}
  \UnaryInf$x : !A\land !B \fCenter !A$
  \RightLabel{\Intro\land}
  \BinaryInf$x : !A\land !B \fCenter !B \land !A$
  \RightLabel{\Intro\lif}
  \UnaryInf$\fCenter (!A \land !B) \lif (!B \land !A)$
  \DisplayProof
  \]
  We assign proof terms from the top down. The proof terms for the
  axioms are just the same as the variable in the context,~$x$. The
  proof terms associated with $\Elim\land[i]$ are then determined as
  $\proj{2}{x}$ and~$\proj{1}{x}$. When we apply \Intro{\land}, we
  combine these two proof terms: $\pair{\proj{2}{x}}{\proj{1}{x}}$.
  Finally, the \Intro{\lif} rule applies a $\lambd$-abstraction,
  binding~$x$ and recording that $x$ labelled the assumption~$!A \land
  !B$: $\lambd[\typeof{x}{!A \land
  !B}][\pair{\proj{2}{x}}{\proj{1}{x}}]$.
  The !!{proof} in \Log{tN2i} then is:
  \[
  \Axiom$x : !A\land !B \fCenter x: !A\land !B$
  \RightLabel{\Elim\land[2]}
  \UnaryInf$x : !A\land !B \fCenter \proj{2}{x} : !B$
  \Axiom$x : !A\land !B \fCenter x: !A\land !B$
  \RightLabel{\Elim\land[1]}
  \UnaryInf$x : !A\land !B \fCenter \proj{1}{x} : !A$
  \RightLabel{\Intro\land}
  \BinaryInf$x : !A\land !B \fCenter \pair{\proj{2}{x}}{\proj{1}{x}} : 
    !B \land !A$
  \RightLabel{\Intro\lif}
  \UnaryInf$\fCenter \lambd[\typeof{x}{!A \land
  !B}][\pair{\proj{2}{x}}{\proj{1}{x}}] : (!A \land !B) \lif (!B \land !A)$
  \DisplayProof
  \]
\end{ex}


\end{document}
